\section{矩阵函数的一阶扰动分析}
\label{sec:multivariable-matrix-perturbation}

上一节求出了矩阵函数的微分。现在考虑另一个问题：已知矩阵$\bm A$和一个小扰动$\bm E$，怎样预测函数值的变化，预测的误差又有多大？微分给出一阶预测，但实际扰动总有有限大小，因此还需要检查近似的余项。

\subsection{微分与局部条件数}

设矩阵函数$\Phi$在$\bm A$处可微。对允许的扰动$\bm E$，其一阶变化为$\mathrm d\Phi_{\bm A}(\bm E)$。这是关于$\bm E$的线性映射；固定方向时，它描述该方向上的变化，并不代表其他方向有同样的响应。

若输入和输出空间分别选定范数，微分的算子范数
\[
 \lVert\mathrm d\Phi_{\bm A}\rVert
 =\sup_{\lVert\bm E\rVert=1}\lVert\mathrm d\Phi_{\bm A}(\bm E)\rVert
\]
给出单位扰动的一阶响应上界。如果只允许对称扰动等具有特定结构的变化，取上确界时也只能使用允许的方向。

输入和输出的尺度均非零时，可将这个上界乘以$\lVert\bm A\rVert/\lVert\Phi(\bm A)\rVert$，衡量相对扰动的一阶放大程度。它描述的是数学问题对输入变化的敏感性，不能单独说明某个求解算法是否稳定。

可微性只保证余项是$o(\lVert\bm E\rVert)$，不会自动给出某个有限扰动的数值误差界。若连接$\bm A$和$\bm A+\bm E$的线段位于定义域中，$\Phi$在该线段的一个开邻域内连续可微，并且微分在该线段上满足$\lVert\mathrm d\Phi_{\bm A+t\bm E}-\mathrm d\Phi_{\bm A}\rVert\leq Lt\lVert\bm E\rVert$，则沿线段应用一元微积分基本定理，得到
\[
 \left\lVert\Phi(\bm A+\bm E)-\Phi(\bm A)-\mathrm d\Phi_{\bm A}(\bm E)\right\rVert
 \leq\int_0^1 Lt\lVert\bm E\rVert^2\,\mathrm dt
 =\frac L2\lVert\bm E\rVert^2.
\]
这个二次上界依赖所给的邻域和常数$L$；下列例子分别核对这些条件。

\subsection{求逆的一阶近似}

取可逆方阵$\bm A$，由式~\eqref{eq:analysis-inverse-differential}，一阶预测为
\[
 (\bm A+\bm E)^{-1}\approx\bm A^{-1}-\bm A^{-1}\bm E\bm A^{-1}.
\]
使用向量范数诱导的矩阵范数（其单位矩阵范数为$1$，并满足乘积不等式），令$q=\lVert\bm A^{-1}\rVert\lVert\bm E\rVert<1$。此时$\bm I+\bm A^{-1}\bm E$可逆，其逆的范数不超过$1/(1-q)$。将$(\bm I+\bm A^{-1}\bm E)^{-1}$展开到一次项，并保留余项，得到
\[
 \left\lVert(\bm A+\bm E)^{-1}-\bm A^{-1}
 +\bm A^{-1}\bm E\bm A^{-1}\right\rVert
 \leq\frac{\lVert\bm A^{-1}\rVert^3\lVert\bm E\rVert^2}{1-q}.
\]
因此，扰动足够小时，预测误差按二次量缩小；当$q$接近$1$时，此界也提醒我们不能把一阶预测当成精确结果。线性方程中同时改变系数和右端项时，还需分别分析它们对解的影响。

\subsection{行列式的一阶近似}

在可逆矩阵$\bm A$处，行列式微分给出
\[
 \det(\bm A+\bm E)=\det\bm A\bigl(1+\operatorname{tr}(\bm A^{-1}\bm E)\bigr)+O(\lVert\bm E\rVert_F^2).
\]
其中迹是相对于原行列式的一阶变化；它随扰动方向而变。在固定有界邻域中，行列式作为矩阵元素的多项式，其二阶导数有界，所以可由沿$\bm A+t\bm E$的Taylor公式得到二次余项界。

奇异矩阵处仍可分析行列式的绝对变化，因为行列式是多项式；但此时$\det\bm A=0$，上式的相对变化和$\bm A^{-1}$都没有定义。行列式的数值小也不必然意味着求逆病态：例如$\bm A=\varepsilon\bm I$（$\varepsilon>0$）的行列式为$\varepsilon^d$，而它的二范数条件数仍为$1$。

\subsection{对数行列式的一阶近似}

设$\det\bm A>0$。取足够小的扰动，使整条线段都留在正行列式区域；由式~\eqref{eq:analysis-logdet-gradient}，有
\[
 \log\det(\bm A+\bm E)-\log\det\bm A
 =\operatorname{tr}(\bm A^{-1}\bm E)+o(\lVert\bm E\rVert_F).
\]
若原来的对数行列式接近零，以它为分母的相对误差不适合描述变化，此时应直接比较绝对变化及其尺度。

对称正定情形可以进一步给出具体余项界。设$\bm A$的最小特征值为$m>0$，$\bm E$对称且$\lVert\bm E\rVert_2<m$。对任意单位向量$\bm x$和$0\leq t\leq1$，
$\bm x^\top(\bm A+t\bm E)\bm x\geq m-\lVert\bm E\rVert_2>0$，因此整条线段保持正定。令$g(t)=\log\det(\bm A+t\bm E)$，则
\[
 g''(t)=-\operatorname{tr}\!\left((\bm A+t\bm E)^{-1}\bm E
 (\bm A+t\bm E)^{-1}\bm E\right).
\]
逆矩阵的二范数至多为$(m-\lVert\bm E\rVert_2)^{-1}$，由Frobenius内积和乘积估计，
$|g''(t)|\leq\lVert\bm E\rVert_F^2/(m-\lVert\bm E\rVert_2)^2$。一元Taylor公式于是给出
\[
 \left|\log\det(\bm A+\bm E)-\log\det\bm A
 -\operatorname{tr}(\bm A^{-1}\bm E)\right|
 \leq\frac{\lVert\bm E\rVert_F^2}{2(m-\lVert\bm E\rVert_2)^2}.
\]

例如$\bm A=\operatorname{diag}(2,1)$，$\bm E=\operatorname{diag}(0.1,0)$时，一阶预测为$0.05$，实际变化为$\log(1.05)\approx0.04879$。此时$m=1$，上述上界约为$0.00618$，大于实际误差约$0.00121$。这个界保证误差不超过多少，并不声称预测误差必然接近上界。
